The prior is expressed through \texttt{tld\_len}, which is a column of the hand-built schema, so we asked whether a representation without it inherits the blind spot. The character-CNN of Section~\ref{sec:data}, which reads the raw URL string, scored on the same phishing-temporal test rows and seeds and read at the same $\tau = 0.5$, still misses $0.771$ of bare \texttt{.vn} phishing and $0.523$ of the other two-character ccTLDs, against $0.074$ on suffixes of three characters or more (5 seeds). The harness fits its CatBoost reference on the whole training window rather than on the oldest $85\%$ of it that the calibrated runs of Table~\ref{tab:suffix} use, and that reference gives $0.632$ and $0.021$ on the same two strata against the table's $0.619$ and $0.022$. The blind spot is therefore not an artefact of the feature schema: it survives a model that never sees a suffix length. What the representation buys is a smaller version of it, pooled \texttt{.vn} falling from $0.899$ to $0.709$, bought by tripling the miss rate on the long-suffix stratum that holds most of the phishing, so it redistributes the error rather than repairing it, which is the budget the rest of this section prices.%
